\Needspace{12\baselineskip}
\originalchapter{补充原题与公开数据分析}
本章根据完整性复核增加七道实际公开原题。数据来自MIT作业页面：基因组FASTA、Seattle SNP项目频率文件、PurR位点和基因表。原始数据文件保持不变，输入哈希、随机背景种子及全部输出见源码包；它们是已有公开数据的重新分析，不是本项目开展的新实验。解答仍为编者独立推导。

\originalsection{PS3.2：同源与异源二聚体}
\begin{exercise}
蛋白接触面为恰含 $N$ 个氨基酸的环，序列随机选取。两异源环 $a,b$ 在相对移位 $s$ 下能量 $E_s^{ab}=\sum_{i=1}^NV(a_i,b_{i+s})$，结合能为全部移位中的最小值。未经翻转的同源环给 $E_s^{aa}=\sum_iV(a_i,a_{i+s})$，物理上将一个环翻转后给 $E_s^{aaR}=\sum_iV(a_{N-i},a_{i+s})$（指标按环模 $N$ 取）。假定大 $N$ 下 $E_s$ 是高斯变量，接触均值零、方差 $\sigma^2$。(a)求异源最小能的均值和分布；(b)求未经翻转同源的分布和均值，注意不同移位实现的数目；(c)求翻转同源的对应量，注意重复接触数；(d)随机选取时哪组最好/最差，若特别设计序列，同源是否必有优势？
\end{exercise}
\begin{solution}
题面给单个接触的两个矩，并不能严格推出共享氨基酸的不同接触及不同移位彼此独立。以下说明原题所意图的极值近似：除完全重复或对称的接触外，将接触看成独立，$V(a,b)=V(b,a)$，不同非等价移位能量近似独立。若不添加这些条件，联合分布和极值均值一般不由两个矩唯一确定。

(a)异源每个移位方差 $N\sigma^2$，近似 $N$ 个独立高斯；其累计最小值为 $1-[1-F(E)]^N$，中心 $-\sigma\sqrt{2N\ln N}$，宽度约 $\sigma\sqrt{N/(2\ln N)}$，是反向Gumbel。

(b)对未经翻转同源，$s$ 和 $N-s$ 的接触集完全相同。排除 $s=0$ 的逐位自接触后，通常只有约 $N/2$ 个不同的能量，方差仍约 $N\sigma^2$。故最小能中心约 $-\sigma\sqrt{2N\ln(N/2)}$，分布为相同高斯最小值式但样本数换为约 $N/2$。偶数 $N$ 的 $s=N/2$ 中每个接触重复两次，单个特殊移位方差约 $2N\sigma^2$；只有一个这样的变量，不改变大 $N$ 的主导极值。$s=0$ 的 $V(a,a)$ 分布并未单独给定，若将它纳入，需另给其统计，不能暗中假设它与异源完全相同。

(c)翻转接触是一个配对对合，除 $O(1)$ 个可能固定位置外，每种无序接触出现两次。因此每个能量约有 $N/2$ 个独立项，每项系数2，方差约 $2N\sigma^2$。约 $N$ 个非等价移位给中心 $-\sigma\sqrt{4N\ln N}$，宽度约 $\sigma\sqrt{N/\ln N}$，相应反向Gumbel将标准差换为 $\sqrt{2N}\sigma$。这个优势来自接触重复的方差增强，而不是平均接触能变负。

(d)在以上随机近似下翻转同源最优，未经翻转同源稍差于异源；后两者主导阶相同，差异属于对数样本数修正。专门设计不是随机抽样的同一问题：若独立优化每个接触，两组都可偏好低能配对；同源减少序列自由度，也可能增加几何匹配限制。仅由随机模型不能证明设计后的同源必有优势，需给定可设计序列、接触势及约束。
\end{solution}

\originalsection{PS4.3：噬菌体包装的能量尺度}
\begin{exercise}
$\phi29$ 噬菌体将20000碱基对的dsDNA装进题面给的圆柱衣壳，半径 $r=42\,\mathrm{nm}$、高 $h=47\,\mathrm{nm}$。DNA由外向内形成螺旋壳，相邻链约 $2.3\,\mathrm{nm}$，室温300 K，公开题面给包装功约 $10^5\kb T$。(a)取持久长约50 nm，估计溶液DNA的构象熵，熵损是否足够；(b)以 $\kappa=\kb T\ell_p$ 估计螺旋弯曲能；(c)电荷间距 $b=.17$ nm、屏蔽长1 nm，邻链距离约2 nm，估计Debye--H\"uckel排斥能，能否解释包装功？
\end{exercise}
\begin{solution}
DNA轮廓长 $L=20000(.34\,\mathrm{nm})=6800$ nm，约有 $L/(2\ell_p)=68$ 个Kuhn节段。每节段方向的构象熵依赖角分辨与约束，绝对连续熵没有无参照的数值；但每段损失 $O(\kb)$ 的独立构象选择给 $O(10^2\kb T)$ 的功尺度，远小于 $10^5\kb T$。它不能独自解释包装功。

若六角紧排的每链截面积 $A_c=\sqrt3(2.3)^2/2\simeq4.58\,\mathrm{nm^2}$，以题给圆柱和同心壳近似有 $\pi h(r^2-r_i^2)=A_cL$，得 $r_i\simeq39.4$ nm。环状壳的弯曲能
\[
\frac{E_b}{\kb T}=\frac{\ell_p}2\int\frac{\dd L}{R^2}
=\frac{\pi\ell_ph}{A_c}\ln\frac r{r_i}\simeq10^2.
\]
直接用 $\ell_pL/(2r^2)$ 也给相同数量级。题面的半径使包装仅占较薄外壳；若半径/直径或真实衣壳形状不同，此估计会变。这里按原题数值作答，没有修改实物尺寸来追求预设结果。

对于两平行线，电荷线密度为 $1/b$，距 $d$，连续线积分给单位长度相互作用
\[
\frac{U_{12}}{\kb TL}=\frac{\ell_B}{b^2}\int_{-\infty}^{\infty}
\frac{\e^{-\sqrt{z^2+d^2}/\lambda}}{\sqrt{z^2+d^2}}\dd z
=\frac{2\ell_B}{b^2}K_0(d/\lambda).
\]
六个邻链每对只数一次，总能约 $3L\,[2\ell_BK_0(d/\lambda)/b^2]$。$\ell_B=.7$ nm、$d=2$ nm、$K_0(2)\simeq.114$ 给 $E_{\rm el}\simeq1.1\times10^5\kb T$，$d=2.3$ nm 则约 $8\times10^4\kb T$，所以裸电荷屏蔽模型能达到所需尺度。

但此结果不是高密度衣壳的精确自由能：凝聚会降低有效电荷，非线性屏蔽、短程水合和离子关联也可显著改变能量。若简单将有效线电荷减为Manning值，排斥按电荷密度平方下降。结论是裸静电具有足够数量级，而不是已唯一证明全部包装功来自Debye势。
\end{solution}

\originalsection{PS6.2：受力DNA模型的完整相界}
\begin{exercise}
将Poland--Scheraga模型推广到两端受力 $F$ 的DNA。(a)对各段位置积分，证明配分函数为杆段 $R(\ell)$ 与气泡 $B(\ell)$ 的交替乘积，段长和为 $N$；(b)双链段为长 $a\ell$ 的刚杆、结合权重 $w^\ell$，求三维方向积分；(c)气泡由两个长度 $\ell$、共享端点的随机链构成，证明 $B(\ell)=s g^{2\ell}\ell^{-3/2}\exp(\beta^2F^2a^2\ell/12)$；(d)以 $z^N$ 作生成函数，用Bose和 $f_m^+(x)=\sum_{\ell\ge1}x^\ell/\ell^m$ 表示两段变换；(e)求解链临界条件；(f)取 $s=1$，在 $w/g^2$ 与 $\beta Fa$ 平面画相界。
\end{exercise}
\begin{solution}
(a)忽略各段之间的相互作用后，外力做功 $F\sum X_i$ 可逐段分开，长度限制仍保留 $\sum\ell_i=N$；固定端型给 $\widetilde Z=\widetilde R/(1-\widetilde R\widetilde B)$。

(b)令 $h=\beta Fa$，归一化方向平均 $\tfrac12\int_{-1}^1\e^{h\ell u}\dd u=\sinh(h\ell)/(h\ell)$，故 $R=w^\ell\sinh(h\ell)/(h\ell)$，$h\to0$ 恢复 $w^\ell$。

(c)每条高斯链的端点权重为 $g^\ell(3/(2\pi\ell a^2))^{3/2}\exp[-3R^2/(2\ell a^2)]$。两条相乘，积分 $\int\dd^3R\e^{\beta\mathbf F\cdot\mathbf R-3R^2/(\ell a^2)}$ 后高斯配平方给 $\e^{h^2\ell/12}$，其余长度前因子为 $\ell^{-3/2}$，微观常数归入 $s$。

(d)定义 $x=g^2z\e^{h^2/12}$，则 $\widetilde B=s f_{3/2}^+(x)$，
\[
\widetilde R=\frac{f_1^+(wz\e^h)-f_1^+(wz\e^{-h})}{2h}
=\frac1{2h}\ln\frac{1-wz\e^{-h}}{1-wz\e^h}.
\]
其适用收敛区需 $wz\e^h<1$。当 $h=0$ 取极限 $\widetilde R=wz/(1-wz)$，不要除以零。

(e)气泡分支点为 $z_c=g^{-2}\e^{-h^2/12}$；极点撞击时 $\widetilde R(z_c)\widetilde B(z_c)=1$，故 $\widetilde R(z_c)=1/[s\zeta(3/2)]$。$A=wz_c$，$C=\exp[2h/(s\zeta(3/2))]$，解得 $A=(C-1)/(C\e^h-\e^{-h})$。

(f)相界为
\[
\left(\frac w{g^2}\right)_c
=\e^{h^2/12}\frac{\exp[2h/(s\zeta(3/2))]-1}
{\exp[h+2h/(s\zeta(3/2))]-\e^{-h}},\quad
\left(\frac w{g^2}\right)_c\bigg|_{h=0}=\frac1{1+s\zeta(3/2)}.
\]
较大结合权重位于双链分支；较小权重位于解链分支。下图用 $s=1$ 绘制模型相界。大 $h$ 时高斯单链权重含 $h^2$，而刚杆含 $h$，造成不应直接用于真实有限伸长DNA的远强力行为。完整模型图与真实物质相图必须区分。
\begin{center}\begin{tikzpicture}\begin{axis}[width=.72\linewidth,height=5cm,xlabel={$h=\beta Fa$},ylabel={$(w/g^2)_c$},xmin=0,xmax=6,grid=major]\addplot[main,thick] table[x=h,y=w_over_g2_critical,col sep=comma]{data/dna-force-boundary.csv};\end{axis}\end{tikzpicture}\end{center}
\end{solution}

\originalsection{PS1.2与PS1.4：公开基因组的ORF和相关}
\begin{exercise}
（PS1.2，可选）使用课程提供的U00096.fna大肠杆菌完整序列。(a)对正向三种阅读框，用5'到3'方向每次读取不重叠三联体，记录连续终止密码间的距离 $L$，可略逆向链；(b)画ORF长度分布，与随机序列比较；(c)估计实际长ORF数量显著超出随机期望的长度阈值。
\end{exercise}
\begin{solution}
实际输入长4639221个碱基。将 $L$ 定义为两终止间\emph{非终止三联体数}，因此位置差需减1；若原始距离按含终止计数，横轴平移一而不改变长尾结论。三正向框共有179558个完全观测区间。每个框首个终止前与最后终止后的截断区间不计入，三框未跨闭环拼接；相对全基因组这只影响至多六个边界区间。

等频独立模型的 $P(L=\ell)=p_s(1-p_s)^\ell$、尾概率 $P(L\ge\ell)=(1-p_s)^\ell$。为了与实际完整区间数比较，将其乘179558；另以实测碱基频率计算独立背景 $p_s\simeq.04562$。下图显示累计尾，避免极长区间稀疏分箱掩盖信号。
\begin{center}\begin{tikzpicture}\begin{semilogyaxis}[width=.78\linewidth,height=5cm,xlabel={非终止三联体数 $L$},ylabel={区间数或期望数},xmin=0,xmax=900,ymin=1e-12,ymax=2e5,legend style={font=\small}]
\addplot[main,thick] table[x=length,y=observed,col sep=comma]{data/orf-tail.csv};\addlegendentry{公开序列}
\addplot[blue,thick] table[x=length,y=composition_background,col sep=comma]{data/orf-tail.csv};\addlegendentry{同组成独立背景}
\end{semilogyaxis}\end{tikzpicture}\end{center}
以600三联体为一个保守阈值，实际有176个区间，同组成独立背景预期仅 $1.22\times10^{-7}$，偏离极显著。这个 $L_{\rm cut}=600$ 不是从唯一统计检验优化的最小阈值，而是明确可复核、足以显示长尾的保守选择；表中其他阈值亦可比较。随机不显著不等于没有功能，随机显著也不等于已证明基因身份，需起始密码、方向和注释等信息。
\end{solution}

\begin{exercise}
（PS1.4，可选）独立位置突变模型预测 $p_{XY}=p_Xp_Y$。(a)计算课程E.coli序列的四碱基频率；(b)计邻位16种碱基对，求 $c_{XY}^{(1)}=p_{XY}^{(1)}/(p_Xp_Y)$；(c)对更远间距重复并讨论相关随距离如何变化。
\end{exercise}
\begin{solution}
按A、C、G、T排序，频率分别为 $.24619,.25423,.25366,.24592$。闭环序列以模基因组长度配对，邻位矩阵为
\[
c^{(1)}\simeq\begin{pmatrix}
1.201&.884&.821&1.103\\1.120&.906&1.159&.814\\.922&1.283&.905&.883\\.755&.921&1.113&1.210
\end{pmatrix}.
\]
例如GT、CG与TA等组合并非全部接近1，独立模型受到实际数据反驳。由于同一序列各对重叠，不能把所有二联体当成完全独立试验后直接套最简单的标准误。

用互信息 $I_n=\sum p_{XY}^{(n)}\log_2c_{XY}^{(n)}$ 汇总，$n=1,2,3,6,9,30,100,300,1000$ 时依次约 \[.01841,\ .00746,\ .00713,\ .00131,\ .00155,\ .00154,\ .00025,\ .00070,\ .000008\] 比特。总体减小但不单调，存在阅读框/长程组成等结构。几次距离的数值不足以确立一个统一指数衰减定律，源码保留各16元矩阵以便检查。
\end{solution}

\originalsection{PS1.6：公开SNP数据与不可识别的拟合}
\begin{exercise}
课程提供Seattle SNP项目的cSNPsAfricanEuropean.dat，每行AD、ED两个群体等位频率。(a)把 $x>.5$ 换成 $1-x$，画频率直方图并与课堂稳态比较，讨论 $N\mu$ 与拟合；(b)比较两群体，取 $\mu=10^{-8}$，用约20--30格、忽略极稀有变体尝试估计有效种群数；(c)用Type.dat比较同义/非同义变体的频率、模型适配及哪类更多，解释可能机制。
\end{exercise}
\begin{solution}
实际原始文件37164行。先剔除在某群体中 $x=0,1$ 的非多态记录，再折叠到 $(0,.5)$；恰为 $.5$ 的边界记录独立处理，拟合未收入它们。用25格、格宽 $.02$ 画图，并按原题忽略最低格后在 $(.02,.5)$ 的24格上拟合。对本卷单倍体拷贝数约定，理论零选择对称突变分布为 $p(x)\propto[x(1-x)]^{\alpha-1}$，$\alpha=2N\mu$。折叠和截断后，每格概率用Beta分布积分再按保留区间归一化，而不是用格中心密度当精确概率。
\begin{center}\begin{tikzpicture}\begin{axis}[width=.76\linewidth,height=4.8cm,xlabel={折叠频率 $x$},ylabel={公开记录数},legend style={font=\small},xmin=0,xmax=.5]
\addplot[main,thick] table[x=x,y=AD,col sep=comma]{data/snp-histogram.csv};\addlegendentry{AD}
\addplot[blue,thick] table[x=x,y=ED,col sep=comma]{data/snp-histogram.csv};\addlegendentry{ED}
\end{axis}\end{tikzpicture}\end{center}
AD、ED分别保留19155、25782条用于拟合。条件多项似然在搜索区间 $\alpha\in[10^{-4},3]$ 均将最优推到\emph{下边界}，而Pearson形状偏差约1252与2174（24格）；采集选择与位点相关使它们不能直接作为校准后的显著性检验量。$\alpha=10^{-4}$ 代入会形式上给 $N=5000$，但这是人为搜索下界而非数据识别出来的种群数。本卷因此不把5000包装成有效种群估计，也不据两个同样贴边的拟合强断言群体大小相同。

这个失败本身回答原题：当前被筛选、四舍五入、带群体结构且不包含全部不变位点的数据，不能在所采用的平稳单参数模型下可靠估计 $N$。改变最低频率截断或模型会改变结果；人口历史与选样过程需要额外信息。图中群体分布确有差别，但不能唯一归因于 $N$。

Type文件给非编码/UTR35536条、非同义842条、同义786条。AD中非同义与同义的多态记录427、406，折叠频率均值分别 $.1322,.1558$；ED对应809、768与 $.0870,.1088$，这里仍按 $0<x<.5$ 口径。非同义变体平均更低与净化选择相容，但机会位点数量、密码简并和发现偏差未知，842比786不能解释成单位位点非同义突变率更高。分组原始计数和条件均值均在QA中，未宣称已完成现代人口遗传推断。
\end{solution}

\originalsection{PS7.1：权重矩阵和全基因组候选}
\begin{exercise}
给PurR结合序列purR\_sites.txt。(a)建每位置每碱基权重并算信息量；(b)用随机DNA输入权重矩阵画得分直方图，取 $T_*=T$ 将其转换为结合能分布；(c)阈值略高于训练位点平均结合能，扫描课程E.coli基因组，检查训练序列是否找全、改变阈值并讨论假阳性，可参考gene\_table.txt；(d)简化模型中长 $L$ 共识序列每失配增加 $\epsilon$，求信息量、随机命中共识的概率及两者联系。
\end{exercise}
\begin{solution}
(a)输入35条、各16碱基。每位置对每种碱基加伪计数 $.5$，所以 $p_i(b)=(n_i(b)+.5)/37$；以实测基因组频率 $q_b$ 为背景，取 $w_i(b)=\ln[p_i(b)/q_b]$。信息量 $I=\sum p_i\log_2(p_i/q_b)\simeq15.866$ 比特。伪计数和背景必须同时报告；零频率导致无限能量的做法会把训练未见字母全部拒绝。

(b)用固定种子20261012生成100000条独立同组成随机16碱基\emph{背景模拟}，$S=\sum w_i$、$E/(\kb T)=-S$，能量零点按此定义。随机平均得分约 $-15.663$，训练平均 $12.519$。这是一组随机序列模拟，不是实测结合能；$T_*=T$ 也是原题的额外假设。

(c)两链全扫描，取 $E_{\rm cut}=\avg E_{\rm train}+\delta\kb T$。$\delta=-1,0,.5,1$ 时阈值得分为13.519、12.519、12.019、11.519，候选窗口数为22、36、39、44；原始35条均能在该基因组的一条方向找到，但相应阈值只收入12、19、19、24条训练记录。故平均能量稍上方的阈值并不能找全输入，原因是低亲和训练位点也是真实训练数据的一部分。阈值下降降低假阴性，同时增加未验证候选。

“额外候选”并不等于已经证实的假阳性：原题只给部分阳性位点，没有完整阴性真值。本卷公开候选的1起始坐标、链方向与得分，保留注释表作核查。注释表标题长4639675与本FASTA长4639221不完全一致，表明版本/坐标必须先核对，不能把靠近某个注释基因就当作已验证调控关系。这里因此不给无证据的精确假阳性数和功能名称。

(d)均匀独立背景中唯一共识命中概率 $4^{-L}$，完全指定共识的信息量 $I=2L$，所以 $P=2^{-I}$。若每个位置失配代价 $\epsilon$ 并按Boltzmann选择，则共识字母概率 $p_c=1/[1+3\e^{-\beta\epsilon}]$，每个非共识字母概率 $p_m=\e^{-\beta\epsilon}/[1+3\e^{-\beta\epsilon}]$，信息量 $I=L[2+p_c\log_2p_c+3p_m\log_2p_m]$。有限容忍时信息量小于 $2L$，\emph{精确共识}的背景概率仍是 $4^{-L}$；$2^{-I}$ 估计的是具有相应典型限制的序列族规模，二者不能不加区分地相等。
\end{solution}

\Needspace{12\baselineskip}
\originalchapter{对齐计数、接触结构与受力模型}
\originalsection{PS2.3：比对数量的约定与题面缺口}
\begin{exercise}
原作业将本题归于Durbin等教材第二章。本卷用独立数学表述保留问题：两有序串分别长 $n,m$，求保持各自内部次序、交错排列成 $n+m$ 串的数量；再讨论把比对逐列上、下符号交错读出并删空位，是否能与所有含空位比对一一对应；对两长 $n$ 串用Stirling式估算相应二项系数。教材文字和版面不重新分发。
\end{exercise}
\begin{solution}
选择 $n+m$ 位置中哪 $n$ 个属于第一串，便确定交错串，数量 $\binom{n+m}m$；两个等长串给 $\binom{2n}n\sim4^n/\sqrt{\pi n}$。Stirling推导是 $(2n)!/(n!)^2\simeq\sqrt{4\pi n}(2n/\e)^{2n}/[2\pi n(n/\e)^{2n}]$。

原题(b)的一一对应需要额外规范。按通常允许匹配、插入、删除三种步的全局比对定义，$n=m=1$ 已有三个比对：两个字母同列，以及两种先后相反的分列空位；交错串却只有AB、BA两个。因而删空位映射对普通全体比对不是双射。实际路径数为Delannoy数
\[
D(n,m)=\sum_{k=0}^{\min(n,m)}\frac{(n+m-k)!}{k!(n-k)!(m-k)!},
\]
其中 $k$ 为同时消耗两字母的对角步数。它满足与动态规划同型的三项递推，$D(1,1)=3$。若特意定义每个交错串的规范配对方式以消去冗余，才得到题意所指的二项系数。写清约定既保留原题组合计算，也避免把错误的普遍双射当成证明。
\end{solution}

\originalsection{PS5.3：接触矩阵的实际统计和REM}
\begin{exercise}
(a)寻找20种氨基酸相互作用 $U(a,a')$，计算平均及方差，原题建议Miyazawa--Jernigan1996表3；(b)紧密自避免链遍历 $L^3=N$ 的立方晶格全部位置，求非骨架最近邻接触数及大 $N$ 下的比例；(c)紧密步数约 $g^N,g\simeq1.85$，结合接触矩估计随机序列的冻结温度与最低能；(d，可选)选蛋白的公开序列和接触结构，用同一势估天然能及 $E_N/E_c$。
\end{exercise}
\begin{solution}
(a)实际读取原题公开链接的MJ表3，使用\emph{上三角及对角的 $e_{ij}$}，再对称填充；下三角是另一个 $e'_{ij}$，不能直接把整张不对称印刷表当一个相互作用矩阵。对20种字母各 $1/20$、400个有序对等权统计，得到 $\bar U=-3.1660$、$\operatorname{Var}U=2.1711025$，单位分别为表中的RT与RT平方。此处只发表聚合数值，不重分发表格或400个原值；脚本从原公开链接临时读取以复核，保留原作者与哈希。

(b)立方盒最近邻边数为 $3L^2(L-1)$，链骨架占 $N-1$ 条，剩下
\[
n=3L^2(L-1)-(N-1)=2N-3N^{2/3}+1,\qquad n/N\to2.
\]
每个接触只数一次，不能把六个邻居直接乘 $N$。在独立接触REM中，构象能均值为 $n\bar U$、方差为 $n\operatorname{Var}U$。

(c)用 $\ln\Omega=N\ln g$，
\[
E_c=n\bar U-\sqrt{2n\operatorname{Var}(U)N\ln g},\qquad
\kb T_c=\sqrt{\frac{n\operatorname{Var}(U)}{2N\ln g}}.
\]
$n/N\to2$ 时 $E_c/N\simeq-8.64339$ 个表能量单位，$\kb T_c\simeq1.87861$ 个表能量单位。表中RT是接触势的参考能量标度，需先规定参考温度才能换成Kelvin，不能把经验势的“RT单位”当成未知求解温度自身而循环定义。

(d)实际分析课程提供4HHB\_A坐标。采用明示的编者接触定义：同链C$\alpha$距小于8 \AA{}，序号差超过2，每对只数一次，141个C$\alpha$、368接触给 $E_N=-1221.80$，同接触数REM参考 $E_c=-1537.39$，比约 $.7947$。1TEN给89个C$\alpha$、240接触、$E_N=-740.43$、$E_c=-998.71$、比 $.7414$。真实折叠稳定性不由这两个比值证明：接触定义、能量零点、非独立接触、非立方紧排和序列设计都会改变解释。源码保留定义和实际坐标，不把模型能量标成实测热力学能。
\end{solution}

\originalsection{PS5.4：由真实坐标得到Ramachandran图}
\begin{exercise}
用公开蛋白结构求每残基的 $\phi,\psi$ 二面角，原题建议Rasmol并提供1TEN和4HHB\_A。(a)画图比较螺旋、片层位置，可参考PDB的HELIX/SHEET字段；(b)找离开主要角度簇的残基及其氨基酸类型；(c)在结构中定位这些残基，讨论是否靠近配体。原题Rasmol操作的目标由下面独立坐标程序等价完成，不宣称运行了Rasmol。
\end{exercise}
\begin{solution}
定义 $\phi=\angle(C_{i-1},N_i,C_{\alpha i},C_i)$、$\psi=\angle(N_i,C_{\alpha i},C_i,N_{i+1})$。四点的两平面法向用叉积得出，再以atan2给带符号角，避免只用arccos丢掉正负。缺原子、断链或跨链的二面角不计算；程序检查肽键距离。实际得到1TEN的88组、4HHB\_A的139组角。

4HHB\_A的110个HELIX标注点平均约 $(-65.4^\circ,-31.6^\circ)$，其余点分散；1TEN原课程文件没有SHEET字段，因此不伪造其二级结构标注，但66个点位于本文明确选定的负 $\phi$、正 $\psi$ 区域，符合片层为主的几何。这个区域计数不是正式DSSP赋值。
\begin{center}\begin{tikzpicture}\begin{axis}[width=.72\linewidth,height=5.5cm,xlabel={$\phi$（度）},ylabel={$\psi$（度）},xmin=-180,xmax=180,ymin=-180,ymax=180,legend style={font=\small}]
\addplot[only marks,mark=*,mark size=1.2pt,main] table[x=phi,y=psi,col sep=comma]{data/ramachandran-4HHB_A.csv};\addlegendentry{4HHB\_A}
\addplot[only marks,mark=triangle*,mark size=1.2pt,blue] table[x=phi,y=psi,col sep=comma]{data/ramachandran-1TEN.csv};\addlegendentry{1TEN}
\end{axis}\end{tikzpicture}\end{center}
示例用正 $\phi$ 来定位偏离主要负角簇的点，不能当作完整的结构质量判定。1TEN有GLN7、ASP12、GLY37、ASN61、ASP65、GLY77；4HHB\_A有GLY18、GLY51、ARG92。Gly缺大的侧链，允许的正角范围较大，因此正 $\phi$ 的Gly未必是异常；Pro的环约束则缩小自由度。真正异常需残基类型相应的允许区与电子密度/结构质量证据。

4HHB\_A保留43个非水配体原子。上述三点C$\alpha$到最近配体原子的距离约16.49、15.58、8.32 \AA{}，均不在本文以5 \AA{}定义的接触邻域；1TEN课程文件无非水配体。因而这组数据并不支持“离群角总靠近配体”。全部角度和距离保存为CSV/QA，实际数据结论与泛泛的可能机制分开。
\end{solution}

\originalsection{PS6.3：DNA过度伸长的模型比较}
\begin{exercise}
B型DNA每碱基对长度3.4 \AA{}，题给实验在 $65\pm5$ pN附近出现过度伸长。(a)每单位独立取正常3.4 \AA{}或伸长5.8 \AA{}，后者代价 $U$，忽略三维方向，求 $L(F,T)$；(b)按原文引用的Smith--Cui--Bustamante实验图估 $U$ 并比较转变宽度；(c)若相邻单位偏好同态，转变将更尖还是更宽？受限原实验图不重新分发，本卷使用题面给定的中心力与独立模型曲线。
\end{exercise}
\begin{solution}
令 $l_0=.34$ nm、$\Delta l=.24$ nm，单点配分函数 $z=\e^{\beta Fl_0}[1+\e^{-\beta(U-F\Delta l)}]$，所以
\[
L=N\left[l_0+\frac{\Delta l}{1+\e^{\beta(U-F\Delta l)}}\right].
\]
中点满足 $U=F_{1/2}\Delta l$。题给65 pN对应 $U\simeq15.6\,\mathrm{pN\,nm}\simeq3.8\kb T$（室温 $\kb T\simeq4.1\,\mathrm{pN\,nm}$）。独立模型10\%到90\%宽度 $\Delta F_{10-90}=2\kb T\ln9/\Delta l\simeq75$ pN，非常宽，和原题描述的明显窄平台/陡起始不相容；精确实验宽度需要原图或原数据，未编造数字化结果。

相邻同态的正协同增加域壁代价，降低混合段权重，在有限温度一维短程模型中仍无严格无限尖的有限温相变，但能显著收窄交叉区。取二状态变量并加入 $-J\sum s_is_{i+1}$ 可用 $2\times2$ 传递矩阵计算；$J>0$ 使同态域增长。模型也未包含WLC方向、序列异质和DNA断链/解链机制，不能仅靠 $U$ 将过伸长归因于唯一新结构。
\end{solution}

\originalsection{PS2.5与PS6.5：在线任务的题面和真实执行边界}
\begin{exercise}
（PS2.5）使用NCBI BLAST/PSI-BLAST。(a)生成随机氨基酸串，长度10--1500，重复与非冗余库做blastp，检查偶然相似；(b)按题给氨基酸频率生成疏水/非疏水周期串，第一十种归为疏水，解释检索；(c)选真实蛋白，替换 $X=0$ 到100\%残基，比较BLAST与PSI-BLAST的耐受；(d，可选)只在同疏水类别内替换，比较阈值。频率表依CMFILVWYAGTSQNEDHRKP顺序为1.660、2.370、4.100、5.810、9.430、6.580、1.240、3.190、7.580、6.840、5.670、7.130、3.970、4.440、6.360、5.270、2.240、5.160、5.940、4.920（归一化后采样）。本卷不重分发原题所引BLAST论文全文。
\end{exercise}
\begin{solution}
随机库中有高分片段的期望随数据库长度增加，低复杂度/周期性破坏简单独立背景，因而要记录SEG过滤、组成校正、得分矩阵与E值阈值。某条随机查询有局部命中不等于真的同源；不能由任意命中名称声称功能。

本次实际提交到NCBI的nr库：4条随机串（10、50、150、1500），一条每四位前两疏水/后两非疏水的150位周期串，以及原题SH3串的0、20、40、80、100\%替换串；种子20261012，BLOSUM62，SEG开启，组成校正2，E阈值 $.01$，每查询保留10条。查询FASTA、请求RID与原始结果哈希在数据与QA中。另在当前Swiss-Prot库实际完成重复随机、两种替换与三轮PSI-BLAST任务，结果在下一节；今日nr远程请求的状态单独记录。

PSI-BLAST迭代建立位置特异矩阵，能强化保守位点但也可能收入错误邻居导致模型漂移。所谓容忍 $X$ 百分比还依赖替换位置、长度、类别和阈值，没有一个普适常数；需对每个生成串实际检索。对于未跑到的算法或失去结果的请求，本卷给机制和复算方案并标未完成，不编造耐受数值。
\end{solution}


\originalsection{PS2.5实际复算：当前数据库与算法比较}
为完成方法任务，另从UniProt官方站实际下载当前Swiss-Prot审校蛋白序列，含575748条、209017843个残基，以NCBI官方BLAST+2.17.0建立本地库。它与原题的2011数据库及今日nr不同，检索结果是明确的新数据库实验，不能冒作历史结果。原始库和官方二进制不塞进源码包，保存下载URL/哈希、查询串、命令和结果；附带脚本可在有该数据库的环境复算。

(a)对10、50、150、1500位各独立生成3条，共12条随机串，在 $E<.01$、BLOSUM62、SEG开启、组成校正2条件下均无保留命中。(b)采用题给频率、每4位前2疏水后2非疏水的一条150位周期串也无保留命中。这个实际阴性不证明所有随机串或周期串都不会得到假阳性；当前数据库、过滤与组成校正显著影响结果。本次没有虚构一个恰好命中的例子。

(c)原题给SH3串长62，每次从全体其他字母替换恰当数量，0、20、40、60、80、100\%各一次；(d)另按题给疏水两类内替换。BLAST批量运行，PSI-BLAST每条单独运行最多3轮、收入阈值 $.005$，目标数上限10。下表给各方法\emph{返回记录中的最小E值}，PSI包含轮次记录，因此记录总数不能当不同蛋白总数。

任意字母替换\par\begin{center}\begin{tabular}{rrrr}\toprule
替换比例 & BLAST最小E & PSI最小E & BLAST保留目标数\\\midrule
0\% & $4.38\times10^{-36}$ & $4.38\times10^{-36}$ & 10\\
20\% & $5.6\times10^{-26}$ & $5.6\times10^{-26}$ & 10\\
40\% & $8.5\times10^{-12}$ & $5.82\times10^{-28}$ & 7\\
60\% & 0.007 & $2.34\times10^{-20}$ & 4\\
80\% & 无 & 无 & 0\\
100\% & 无 & 无 & 0\\
\bottomrule\end{tabular}\end{center}

同类字母替换\par\begin{center}\begin{tabular}{rrrr}\toprule
替换比例 & BLAST最小E & PSI最小E & BLAST保留目标数\\\midrule
0\% & $4.38\times10^{-36}$ & $4.38\times10^{-36}$ & 10\\
20\% & $2.37\times10^{-24}$ & $2.01\times10^{-25}$ & 10\\
40\% & $3.65\times10^{-17}$ & $4.31\times10^{-25}$ & 10\\
60\% & $2.72\times10^{-5}$ & $4.76\times10^{-28}$ & 5\\
80\% & 无 & 无 & 0\\
100\% & 无 & 无 & 0\\
\bottomrule\end{tabular}\end{center}

本次两机制在60\%替换仍有命中，80\%与100\%均无；它给一个 $60$--$80$\% 的\emph{当前有限测试间隔}，不是通用耐受阈值。同类替换60\%仍保留较高覆盖（最佳BLAST约94\%），任意替换60\%的最佳覆盖约77\%；PSI在已有可靠种子的40/60\%情形明显加强匹配。完全没有合适初始种子的80/100\%情形，迭代未凭空恢复。

“最佳”目标和E值会随算法/轮次改变，不能把任意谱系命中自动赋予蛋白功能。更多突变实现、不同位置、SEG开关和不同数据库应作为进一步实验。源题的主方法已实际运行；2011库无法按本次取源恢复，今日nr远程请求的完成状态另有记录，不用当前Swiss-Prot替代它的身份。

\begin{exercise}
（PS6.5，可选）用原题UCSD或Ohio RNA pulling服务，对长度接近的均匀序列、GC交替、AU交替及实际RNA序列求力延伸曲线，比较通性及与PS6.4理论是否相似。
\end{exercise}
\begin{solution}
该任务必须有服务或等价完整结构/受力计算的真实输出。2026年10月12日，实际网络请求Ohio的HTTP、HTTPS入口均返回503；UCSD入口连接返回403。网页检索缓存仍显示服务说明，但缓存存在不能证明今天可执行。源核验与失败记录保留，此题题面已经收录，原服务数值结果仍是明确缺口。

可合理预期的机制是：全A等不形成Watson--Crick配对的序列接近单链弹性；GC与AU交替允许很多发夹/内部配对，结合能及堆叠不同；天然序列会有特定结构与多个展开阶段。但这些只是机制预期，不是所要求的四条已经测算曲线，也不能由此断言临界力大小。本卷没有用随手画的曲线替代失效服务的输出。
\end{solution}
